% Options for packages loaded elsewhere
\PassOptionsToPackage{unicode}{hyperref}
\PassOptionsToPackage{hyphens}{url}
\PassOptionsToPackage{dvipsnames,svgnames,x11names}{xcolor}
\documentclass[
  10pt,
  fontset=fandol]{ctexart}
\usepackage{xcolor}
\usepackage[margin=16mm]{geometry}
\usepackage{amsmath,amssymb}
\setcounter{secnumdepth}{-\maxdimen} % remove section numbering
\usepackage{iftex}
\ifPDFTeX
  \usepackage[T1]{fontenc}
  \usepackage[utf8]{inputenc}
  \usepackage{textcomp} % provide euro and other symbols
\else % if luatex or xetex
  \usepackage{unicode-math} % this also loads fontspec
  \defaultfontfeatures{Scale=MatchLowercase}
  \defaultfontfeatures[\rmfamily]{Ligatures=TeX,Scale=1}
\fi
\usepackage{lmodern}
\ifPDFTeX\else
  % xetex/luatex font selection
\fi
% Use upquote if available, for straight quotes in verbatim environments
\IfFileExists{upquote.sty}{\usepackage{upquote}}{}
\IfFileExists{microtype.sty}{% use microtype if available
  \usepackage[]{microtype}
  \UseMicrotypeSet[protrusion]{basicmath} % disable protrusion for tt fonts
}{}
\makeatletter
\@ifundefined{KOMAClassName}{% if non-KOMA class
  \IfFileExists{parskip.sty}{%
    \usepackage{parskip}
  }{% else
    \setlength{\parindent}{0pt}
    \setlength{\parskip}{6pt plus 2pt minus 1pt}}
}{% if KOMA class
  \KOMAoptions{parskip=half}}
\makeatother
\setlength{\emergencystretch}{3em} % prevent overfull lines
\providecommand{\tightlist}{%
  \setlength{\itemsep}{0pt}\setlength{\parskip}{0pt}}
\usepackage{amsmath,amssymb,mathtools,needspace}
\xeCJKDeclareCharClass{CJK}{"2032,"0400->"04FF,"2160->"216B,"2460->"2473,"25A0,"25C7,"25CB,"25CF}
\IfFontExistsTF{Arial Unicode MS}{\xeCJKsetup{AutoFallBack=true}\setCJKfallbackfamilyfont{\CJKrmdefault}{Arial Unicode MS}}{}
\setcounter{secnumdepth}{0}
\setlength{\emergencystretch}{2em}
\allowdisplaybreaks
\usepackage{bookmark}
\IfFileExists{xurl.sty}{\usepackage{xurl}}{} % add URL line breaks if available
\urlstyle{same}
\hypersetup{
  pdftitle={函数按正交多项式展开},
  colorlinks=true,
  linkcolor={Maroon},
  filecolor={Maroon},
  citecolor={Blue},
  urlcolor={Blue},
  pdfcreator={LaTeX via pandoc}}

\title{函数按正交多项式展开}
\author{}
\date{}

\begin{document}
\maketitle

\subsection{3.5 函数按正交多项式展开}\label{ux51fdux6570ux6309ux6b63ux4ea4ux591aux9879ux5f0fux5c55ux5f00}

设 \(f(x)\in C[a,b]\)，用正交多项式 \(\{g_0(x),g_1(x),\cdots,g_n(x)\}\) 作基，求最佳平方逼近多项式

\[
S_n(x)=a_0g_0(x)+a_1g_1(x)+\cdots+a_ng_n(x).\tag{3.5.1}
\]

由 \(\{g_k(x)\}\) 的正交性及式 (3.3.13)，可求得系数

\[
a_k=\frac{(f,g_k)}{(g_k,g_k)}\quad(k=0,1,\cdots,n),\tag{3.5.2}
\]

于是，\(f(x)\) 的最佳平方逼近多项式为

\begin{quote}
校注（agent 补充）：式 (3.4.11) 的原扫描明确印为 \(\mathrm e^n\)，已忠实保留；这是疑似原书错误，不是 OCR 误读。取 \(n=0\)，原式给出 \(L_0(x)=\mathrm e^{-x}\)，与紧接着印出的 \(L_0(x)=1\) 矛盾。相容的 Rodrigues 表达式前因子应为 \(\mathrm e^x\)。见 \href{../assets/ch03b/p076-laguerre-formula.png}{原式放大裁切}。
\end{quote}

\href{../../source-images/pdf-077.jpeg}{查看原页}

\[
S_n(x)=\sum_{k=0}^n\frac{(f,g_k)}{(g_k,g_k)}g_k(x).\tag{3.5.3}
\]

由式 (3.3.15) 可得均方误差为

\[
\|\delta_n\|_2=\|f-S_n\|_2=\sqrt{\|f\|_2^2-\sum_{k=0}^n\frac{(f,g_k)}{(g_k,g_k)}(f,g_k)}.\tag{3.5.4}
\]

若 \(f(x)\) 在 \([a,b]\) 上按正交多项式 \(\{g_k(x)\}\) 展开，系数 \(a_k\)（\(k=0,1,\cdots\)）按式 (3.5.2) 计算，这样可得到 \(f(x)\) 的展开式

\[
f(x)\sim\sum_{k=0}^\infty a_kg_k(x).\tag{3.5.5}
\]

式 (3.5.5) 右端级数称为广义 Fourier 级数，系数 \(a_k\) 称为广义 Fourier 系数，它是 Fourier 级数的直接推广。从以上讨论可知，任何 \(f(x)\in C[a,b]\) 均可展开成广义 Fourier 级数，其部分和 \(S_n(x)\) 是 \(f(x)\) 的最佳平方逼近，系数 \(a_k\) 与 \(n\) 无关，因此，当 \(n\) 增加时，只要多计算增加的系数 \(a_k\)。级数 (3.5.5) 可能不一致收敛于 \(f(x)\)，但在 \(f(x)\) 满足一定条件下也可一致收敛到 \(f(x)\)。

下面考虑函数 \(f(x)\in C[-1,1]\) 按 Legendre 多项式 \(\{P_0(x),P_1(x),\cdots\}\) 展开的最佳平方逼近多项式 \(S_n^*(x)\)。由式 (3.5.1) 和式 (3.5.2) 可得

\[
S_n^*(x)=a_0^*P_0(x)+a_1^*P_1(x)+\cdots+a_n^*P_n(x),\tag{3.5.6}
\]

其中

\[
a_k^*=\frac{(f,P_k)}{(P_k,P_k)}=\frac{2k+1}{2}\int_{-1}^1f(x)P_k(x)\,\mathrm dx.\tag{3.5.7}
\]

根据式 (3.5.4)，平方误差为

\[
\|\delta_k\|_2^2=\int_{-1}^1 f^2(x)\,\mathrm dx-\sum_{k=0}^n\frac{2}{2k+1}(a_k^*)^2.\tag{3.5.8}
\]

\textbf{例 3.4} 求 \(f(x)=\mathrm e^x\) 在 \([-1,1]\) 上的三次最佳平方逼近多项式。

\textbf{解} 先计算 \((f,P_k)\)（\(k=0,1,2,3\)），即

\[
(f,P_0)=\int_{-1}^1\mathrm e^x\,\mathrm dx=\mathrm e-\frac1{\mathrm e}\approx2.3504,
\]

\[
(f,P_1)=\int_{-1}^1x\mathrm e^x\,\mathrm dx=2\mathrm e^{-1}\approx0.7358,
\]

\[
(f,P_2)=\int_{-1}^1\left(\frac32x^2-\frac12\right)\mathrm e^x\,\mathrm dx=\mathrm e-\frac7{\mathrm e}\approx0.1431,
\]

\[
(f,P_3)=\int_{-1}^1\left(\frac52x^3-\frac32x\right)\mathrm e^x\,\mathrm dx=37\frac1{\mathrm e}-5\mathrm e\approx0.02013.
\]

由式 (3.5.7) 得

\[
a_0^*=(f,P_0)/2=1.1752,\quad a_1^*=3(f,P_1)/2=1.1036,
\]

\[
a_2^*=5(f,P_2)/2=0.3578,\quad a_3^*=7(f,P_3)/2=0.07046,
\]

代入式 (3.5.6)，得

\begin{quote}
校注（agent 补充）：式 (3.5.8) 左端原书印为 \(\|\delta_k\|_2^2\)，已保留。其右端求和上限为 \(n\)，与前式 (3.5.4) 一致时左端应记为 \(\|\delta_n\|_2^2\)，属疑似原书下标错误。见 \href{../assets/ch03b/p077-error-formula.png}{原式放大裁切}。
\end{quote}

\href{../../source-images/pdf-078.jpeg}{查看原页}

\[
S_3^*(x)=0.9963+0.9979x+0.5367x^2+0.1761x^3.
\]

均方误差为

\[
\|\delta_n\|_2=\|\mathrm e^x-S_3^*(x)\|_2=\sqrt{\int_{-1}^1\mathrm e^{2x}\,\mathrm dx-\sum_{k=0}^3\frac{2}{2k+1}(a_k^*)^2}\leqslant0.0084,
\]

最大误差为

\[
\|\delta_n\|_\infty=\|\mathrm e^x-S_3(x)\|_\infty\leqslant0.0112.
\]

如果 \(f(x)\in C[a,b]\)，求 \(f(x)\) 在区间 \([a,b]\) 上的最佳平方逼近多项式。作变换

\[
x=\frac{b-a}{2}t+\frac{b+a}{2}\quad(-1\leqslant t\leqslant1),
\]

于是

\[
F(t)=f\left(\frac{b-a}{2}t+\frac{b+a}{2}\right).
\]

在 \([-1,1]\) 上可用 Legendre 多项式作最佳平方逼近多项式 \(S_n^*(t)\)，从而得到 \(f(x)\) 在区间 \([a,b]\) 上的最佳平方逼近多项式 \(S_n^*\left(\dfrac1{b-a}(2x-a-b)\right)\)。

由于 Legendre 多项式 \(\{P_k(x)\}\) 是在区间 \([-1,1]\) 上由 \(\{1,\allowbreak x,\allowbreak x^2,\allowbreak \cdots,\allowbreak x^k,\allowbreak \cdots\}\) 正交化得到的，因此利用函数的 Legendre 展开部分和得到最佳平方逼近多项式与由

\[
S^*(x)=a_0+a_1x+\cdots+a_nx^n
\]

直接通过解法方程得到的 \(H_n\) 中的最佳平方逼近多项式是一致的，只是当 \(n\) 较大时求法方程会出现病态方程，计算误差较大，不能使用。而用 Legendre 展开不用解线性方程组，不存在病态问题，计算公式使用起来也较方便，因此通常都用此法求最佳平方逼近多项式。

\begin{center}\rule{0.5\linewidth}{0.5pt}\end{center}

\subsection{相关原页校注（agent 补充）}\label{ux76f8ux5173ux539fux9875ux6821ux6ce8agent-ux8865ux5145}

以下校注来自本节覆盖的原页，可能也涉及同页相邻小节；原文照录与校正说明分开保留。

\subsubsection{原书 PDF 页 78 的校注}\label{ux539fux4e66-pdf-ux9875-78-ux7684ux6821ux6ce8}

\href{../pages/pdf-078.md}{完整原页转录} · \href{../../source-images/pdf-078.jpeg}{原页图像}

校注记录：ch03b-E003。

\begin{quote}
校注（agent 补充）：本页最大误差式原书写 \(S_3(x)\)，而同例的最佳平方逼近多项式记作 \(S_3^*(x)\)；转录保留原来的无星号记法。
\end{quote}

\subsection{本节覆盖原页的校注记录}\label{ux672cux8282ux8986ux76d6ux539fux9875ux7684ux6821ux6ce8ux8bb0ux5f55}

下列记录按原页关联，含同页相邻小节的校注；计算时同时读取上文校注和完整原页。

\begin{itemize}
\tightlist
\item
  \texttt{ch03b-E001}：原书 PDF 页 76
\item
  \texttt{ch03b-E002}：原书 PDF 页 77
\item
  \texttt{ch03b-E003}：原书 PDF 页 78
\end{itemize}

\end{document}
